\begin{table}[H]
\centering
\caption{Clasificación de compuestos inorgánicos}
\label{tab:nomenclatura_compuestos}
\begin{threeparttable}
{\setstretch{1.75}
\begin{tabular}{p{0.2\linewidth}p{0.25\linewidth}p{0.2\linewidth}p{0.25\linewidth}} \hline

    \bfseries{Nomenclatura} & \bfseries{Unión de:} & \bfseries{Tipo de compuesto} & \bfseries{Ejemplo} \\ \hline
    
    Hidrácidos          & $H + \text{ No Metal}$            & Binario      & $HCl$:     \par Ácido clohrídrico \\
    Hidróxidos          & $\text{Metal} + OH^{-}$           & Ternario     & $NaOH$:    \par Hidróxido de sodio \\
    Oxácidos            & $H + \text{No Metal} + O$         & Ternario     & $HNO_3$:   \par Ácido nítrico \\
    Oxisales            & $\text{Metal} + \text{Anión} + O$ & Ternario     & $CaCO_3$:  \par Carbonato de calcio \\
    Óxidos metálicos    & $\text{Metal} + O$                & Binario      & $CaO$:     \par Óxido de calcio \\
    Óxidos NO metálicos & $\text{No Metal} + O$             & Binario      & $H_{2}O$ \\
    Sales binarias      & $\text{Metal} + \text{No metal}$  & Binario      & $NaCl$:    \par Cloruro de sodio \\
    Sales ácidas        & $\text{Metal} + H + \text{Anión}$ & Poli-atómico & $KaHSO_3$: \par Bisulfato de potasio \\
    
    \hline
\end{tabular}
}
\begin{tablenotes} % ex: \tntote{1} ... \item[1]
    \item[] 
\end{tablenotes}
\end{threeparttable}
\end{table}